%! Confidence Intervals for the Difference of Two Means — Unit 7, Lesson 7.6 — Warm-Up
\documentclass[10pt]{article}
\usepackage{apstats-article}
\usepackage{apstats-boxes}
\newcommand{\TallMath}[1]{$\displaystyle #1\rule[-1.4em]{0pt}{3.2em}$}

\begin{document}

\pageheader{Unit 7, Lesson 7.6}{Warm-Up}
\namedateperiod

\vspace{0.15in}

\begin{notesbox}{Spiral Review --- One-Sample $t$-Interval (Lessons 7.2--7.3)}

A school dietitian claims that students consume an average of 850 calories at lunch. A
random sample of 22 students recorded their lunch calories. The results were
$\bar x = 920$ cal and $s = 130$ cal. A dotplot shows no strong skew or outliers.

\textbf{1.} Name the inference method and give the degrees of freedom.

\quad Method: \blank{6cm} \quad $df = $ \blank{1.5cm}

\textbf{2.} Verify both conditions.
\begin{itemize}
  \item Random: \writeline
  \item Normal: \writeline
\end{itemize}

\textbf{3.} Compute the standard error and margin of error for a 95\% confidence interval
($t^* = 2.080$). Show your work.

\vspace{0.1in}
$SE = \dfrac{s}{\sqrt{n}} = \dfrac{\blank{2cm}}{\sqrt{\blank{1.5cm}}} = $ \blank{2.5cm}
\hfill
$ME = t^* \times SE = \blank{2cm} \times \blank{2.5cm} = $ \blank{2.5cm}

\textbf{4.} Write the 95\% confidence interval and interpret it in context.

\quad CI: $\bar x \pm ME = \blank{2cm} \pm \blank{2cm} = ($ \blank{2cm} , \blank{2cm} $)$

\quad Interpretation: \writeline

\end{notesbox}

\end{document}
